\section{Contextualização e Objetivos Acadêmicos}\label{sec:contexto_objetivos}

A modelagem eletromecânica do motor de indução trifásico com rotor em gaiola de esquilo fundamenta-se na representação por circuitos elétricos equivalentes por fase em regime permanente senoidal e no equacionamento rigoroso da dinâmica rotacional \cite{fitzgerald2014maquinas,chapman2013fundamentos}. A predição de grandezas nominais e transitórias decorre da segregação precisa entre as perdas Joule, as perdas mecânicas e a potência eletromecânica convertida, governada pelo escorregamento relativo do campo girante.

% =========================================================================
% [DIRETRIZ DE REDAÇÃO]
% - Objetivo: Contextualizar a máquina de indução como atuador industrial preponderante
%   e formalizar o escopo acadêmico do trabalho final de ELE035 (sim_proj.pdf).
% - Formulação: Definir velocidade síncrona (\omega_s, n_s), escorregamento (s) e a
%   segregação metodológica entre Parte 1 (analítica) e Parte 2 (transitório dinâmico).
% - Figuras/Tabelas: Citar o documento de requisitos acadêmicos e relatórios anteriores.
% - Confronto: Justificar o critério de convergência de regime permanente em <= 2,0%.
% =========================================================================

A velocidade síncrona do campo magnético girante produzido pelos enrolamentos do estator e o escorregamento normalizado $s$ são expressos por:
\begin{equation}
n_s = \frac{120 \cdot f}{p} \quad [\RPM], \qquad \omega_s = \frac{2\pi f}{p_{\text{pares}}} \quad [\rads],
\label{eq:velocidade_sincrona}
\end{equation}
\begin{equation}
s = \frac{n_s - n}{n_s} = \frac{\omega_s - \omega_m}{\omega_s},
\label{eq:escorregamento}
\end{equation}
em que $f$ representa a frequência da rede elétrica, $p$ o número de polos ($p_{\text{pares}} = p/2$), e $n$ e $\omega_m$ a rotação mecânica do rotor.
